\section{Explained: step-by-step walk-throughs}\label{app:explained}
\noindent The following walk-throughs explain, in plain language and step by step, the notions used in the main text (rings, the operator ring, the annihilator, exponential-polynomials, the gcd, factors, the interaction term, the nonlinear case, the algorithm). They come from the explained edition of the first manuscript, where each was a box placed at the point named in its first line. They are self-contained and can be read in any order.

\subsection{Rings, the operator ring, symbol and characteristic variety}\label{expl:E1}
\noindent\emph{(E1: after Figure 1 (Section 2.1))}
\textbf{Step 0: what is a ring?} A \emph{ring} is the analogue of a group in which there are two operations, addition and multiplication. A set $R$ with $+$ and $\cdot$ is a ring if:
\begin{enumerate}[nosep]
\item $(R,+)$ is an abelian group: addition is associative and commutative, there is an additive identity $0$, and every element $a$ has an additive inverse $-a$;
\item multiplication is associative, $(ab)c=a(bc)$;
\item multiplication distributes over addition, $a(b+c)=ab+ac$ and $(a+b)c=ac+bc$.
\end{enumerate}
Unlike in a group, a ring does \emph{not} require every nonzero element to have a multiplicative inverse, and multiplication need not be commutative. The hierarchy is:
\begin{itemize}[nosep]
\item \emph{group}: one operation;
\item \emph{ring}: an additive abelian group together with a multiplication compatible with the addition;
\item \emph{field}: a commutative ring with $1\ne0$ in which every nonzero element has a multiplicative inverse.
\end{itemize}
Examples: $(\Z,+)$ is a group; $(\Z,+,\cdot)$ is a ring but not a field (2 has no inverse in $\Z$); $(\mathbb Q,+,\cdot)$ is a field. A ring is therefore \emph{not} simply a group under multiplication: it is an additive abelian group on which one can also multiply, with multiplication obeying the distributive law.

\medskip\textbf{Step 1: what is $\ring=\C[\partial_t,\partial_x]$?} It is the ring of polynomials in the two formal variables $\partial_t,\partial_x$ with complex coefficients. Elements of $\ring$ are, for instance,
\[
P(\partial)=\partial_t^2+\partial_x^2,\qquad Q(\partial)=3\partial_t\partial_x+2,\qquad S(\partial)=i\partial_t^3-5\partial_x .
\]
Why is it a ring? One can add operators ($P+Q$), multiply them by composing them ($PQ$), there is a zero operator and every operator has an additive inverse, and composition is associative and distributes over addition. The crucial point is that these are \emph{constant-coefficient} operators: the coefficients are complex numbers, not functions of $t$ or $x$. So $2\partial_t^2+3\partial_x+1$ is in $\ring$, whereas $t\,\partial_t$ is not (its coefficient $t$ is not constant).

\medskip\textbf{Step 2: why $\ring\cong\C[\tau,\xi]$?} The symbol $\cong$ means \emph{isomorphic}: the two rings have exactly the same algebraic structure. Here $\tau,\xi$ are ordinary formal polynomial variables and we simply replace
\[
\partial_t\longleftrightarrow\tau,\qquad\partial_x\longleftrightarrow\xi,
\qquad\text{e.g.}\qquad
\underbrace{\partial_t^2+3\partial_t\partial_x+2\partial_x^2}_{\text{differential operator}}\ \longleftrightarrow\ \underbrace{\tau^2+3\tau\xi+2\xi^2}_{\text{polynomial}} .
\]
This is legitimate because constant-coefficient partial derivatives commute, $\partial_t\partial_x=\partial_x\partial_t$, so the operator ring is commutative, exactly like the polynomial ring. At this stage $\tau$ and $\xi$ are \emph{not} numbers: they are formal variables that encode the operators.

\medskip\textbf{Step 3: the symbol.} For $P(\partial)=\partial_t^2-\partial_x^2+1$ the symbol is obtained by replacing each derivative with its variable,
\[
P(\tau,\xi)=\tau^2-\xi^2+1 .
\]
The symbol translates a differential operator into an algebraic object, and polynomial algebra is often easier to study than differential operators directly. \emph{Convention.} Other texts replace $\partial$ by $i\tau,i\xi$ (Fourier convention). In this paper the substitution is the \emph{direct} one, $\partial\mapsto(\tau,\xi)$; the Fourier convention $\partial_x\mapsto i\kappa$ appears only later, in Sections~\ref{sec:groups} and~\ref{sec:inside}, and is stated there.

\medskip\textbf{Step 4: the characteristic variety.} Given the symbol, $\Var(P)=\{(\tau,\xi)\in\C^2:P(\tau,\xi)=0\}$: the set of all pairs of complex numbers at which the symbol vanishes. For $P=\tau^2-\xi^2+1$ it is $\tau^2=\xi^2-1$, a subset of $\C^2$; a pair belongs to $\Var(P)$ exactly when it satisfies this equation. For $P=\tau-\xi$ it is the line $\tau=\xi$ of $\C^2$.

\medskip\textbf{Example: the wave operator.} $P(\partial)=\partial_t^2-c^2\partial_x^2$ has symbol $\tau^2-c^2\xi^2=(\tau-c\xi)(\tau+c\xi)$, so
\[
\Var(P)=\{\tau=c\xi\}\cup\{\tau=-c\xi\},
\]
two lines (Figure~\ref{fig:varieties}, left). They encode the characteristic directions of wave propagation; their precise interpretation depends on the PDE and on whether the covectors are taken real or complex.

\medskip\textbf{The whole passage in one picture.}
\[
\underbrace{P(\partial_t,\partial_x)}_{\text{operator on functions}}\ \xrightarrow{\ \partial\mapsto(\tau,\xi)\ }\ \underbrace{P(\tau,\xi)}_{\text{polynomial}}\ \xrightarrow{\ \text{zeros}\ }\ \underbrace{\Var(P)\subset\C^2}_{\text{set of points}}
\]
Three related but \emph{different} objects: an operator acting on functions, a polynomial, and a set of points of complex space.

\subsection{Annihilators, ideals, (F1)--(F3) and the splitting of the gcd}\label{expl:E2}
\noindent\emph{(E2: after Lemma 2.5 and the paragraph on G\_u, J\_u (Section 2.2))}
\textbf{Step 1: annihilating a function.} Take $u(x)=e^{x}$. Since $\partial_xe^x=e^x$, $(\partial_x-1)e^x=0$: the operator $\partial_x-1$ \emph{annihilates} $e^x$. It is not the only one: $(\partial_x^2-1)e^x=0$, and $\partial_x^2-1=(\partial_x-1)(\partial_x+1)$ is annihilating precisely because it contains the factor $\partial_x-1$. More generally, if $Q(\partial_x)e^x=0$ then every polynomial multiple of $Q$ annihilates $e^x$, for instance $(\partial_x+3)(\partial_x-1)e^x=0$. This is the basic reason annihilators are ideals.

\medskip\textbf{Step 2: the annihilator is the whole ideal $(\partial_x-1)$.} For $R=\C[\partial_x]$, any polynomial can be divided by $z-1$: $Q(z)=(z-1)S(z)+Q(1)$. Replacing $z$ by $\partial_x$ and applying to $e^x$ gives $Q(\partial_x)e^x=Q(1)e^x$, hence
\[
Q(\partial_x)e^x=0\iff Q(1)=0\iff(z-1)\mid Q(z).
\]
So $\Ann(e^x)=(\partial_x-1)=\{S(\partial_x)(\partial_x-1):S\in\C[\partial_x]\}$, and every annihilating polynomial is a multiple of $\partial_x-1$.

\medskip\textbf{Step 3: why an ideal.} An ideal is a subset closed under addition (and additive inverses) and under multiplication by any element of the ring. For annihilators: if $Qu=0$ and $Pu=0$ then $(Q+P)u=Qu+Pu=0$; if $Qu=0$ and $S\in\ring$ then $(SQ)u=S(Qu)=S(0)=0$. Intuition: once an operator kills a function, composing it with another constant-coefficient operator cannot bring the function back to life.

\medskip\textbf{Step 4: one variable versus two variables.} In $\C[\xi]$ every ideal is generated by a single polynomial, $I=(g)$ -- like the multiples of $6$ in $\Z$. In $\C[\tau,\xi]$ several generators may be needed. Consider $I=(\tau,\xi-1)=\{\tau A+(\xi-1)B\}$: neither generator is a multiple of the other and the ideal is not generated by one polynomial. This matters because the annihilator of a function of two variables can have several independent generators: for $u(t,x)=e^x$ we have $\partial_tu=0$ and $(\partial_x-1)u=0$, and
\[
\Ann(u)=(\partial_t,\ \partial_x-1).
\]
Both equations are needed to describe it in the two-variable ring.

\medskip\textbf{Step 5: the three facts.}
\begin{itemize}[nosep]
\item \textbf{(F1) Unique factorization.} Every nonzero polynomial factors into irreducibles, uniquely up to constants and order; e.g.\ $\tau^2-\xi^2=(\tau-\xi)(\tau+\xi)$. This allows one to define the gcd of an ideal: the greatest common polynomial factor of all its elements.
\item \textbf{(F2) B\'ezout.} Two polynomials in $(\tau,\xi)$ with no common nonconstant factor have finitely many common zeros in $\C^2$. Example: $\tau$ and $\xi-1$ have no common factor and their only common zero is $(0,1)$. Geometrically, two algebraic curves without a shared component meet in finitely many points.
\item \textbf{(F3) Nullstellensatz.} If $q$ vanishes on $\Var(I)$ then $q^k\in I$ for some $k\ge1$: $q|_{\Var(I)}=0\Rightarrow q^k\in I$. It connects the geometry of a zero set to the algebra of the ideal. The power matters: $q$ itself need not belong to $I$ (e.g.\ $I=(\tau^2)$, $q=\tau$).
\end{itemize}

\medskip\textbf{Step 6: the splitting lemma.} Every nonzero ideal $I\subset\ring$ can be written $I=gJ$, with $g=\Gcd(I)$ and $\Var(J)$ finite. Here $g$ is the polynomial factor shared by every element of $I$; $J$ is the ideal of the quotients $f/g$; because $g$ is the gcd, the elements of $J$ have no common nonconstant factor, and by (F2) their common zero set is finite. So the ideal splits into a common polynomial factor and a residual ideal with finite zero set.
\begin{itemize}[nosep]
\item \emph{Example A (shared factor).} $I=(\tau(\xi-1),\tau(\xi+1))$. Both generators contain $\tau$, so $g=\tau$ and $J=(\xi-1,\xi+1)$. Since $(\xi+1)-(\xi-1)=2$, $J=(1)=\ring$, whose zero set is empty (hence finite). Thus $I=(\tau)$.
\item \emph{Example B (no shared factor).} $I=(\tau,\xi-1)$: $g=1$, $J=I$, $\Var(J)=\{(0,1)\}$, finite.
\end{itemize}

\medskip\textbf{Step 7: applied to $u$.} Define $G_u=\Gcd\Ann(u)$ and call the remaining ideal $J_u$, so $\Ann(u)=G_uJ_u$. For $u=e^x$ in $\C[\partial_t,\partial_x]$ we found $\Ann(e^x)=(\partial_t,\partial_x-1)$; the generators have no common nonconstant factor, so
\[
G_u=1,\qquad J_u=(\partial_t,\partial_x-1),\qquad \Var(J_u)=\{(0,1)\},
\]
the single point corresponding to the exponential $e^{0\cdot t+1\cdot x}=e^x$. A point of the characteristic picture corresponds to a possible exponential mode. \emph{Qualification:} the reading of $G_u$ as the ``genuinely infinite-dimensional part'' and of $J_u$ as the ``finitely many exponentials present'' is an intuition for the later theory, not a universal identification for arbitrary functions without further hypotheses.

\subsection{The finite dimension criterion for exponential-polynomials}\label{expl:E3}
\noindent\emph{(E3: after Lemma 2.7 (EP), before the Takeaway of Section 2.3)}
The lemma says: \emph{a function is a finite sum of polynomial-times-exponential terms if and only if all its derivatives span a finite-dimensional vector space.}

\medskip\textbf{Step 1: what is $\EP$?} The set of functions $w(t,x)=\sum_{j=1}^Np_j(t,x)e^{\tau_jt+\xi_jx}$, with $p_j$ polynomials and $\tau_j,\xi_j\in\C$. Examples: $e^{2t+x}+3e^{-t+4x}$ and $t^2e^x+(t+x)e^{t-x}$. The second has polynomial prefactors, so $\EP$ is more general than sums of pure exponentials. Polynomial prefactors appear because differentiating such a function gives another function of the same form: $\partial_x(te^{2t+x})=te^{2t+x}$ and $\partial_t(te^{2t+x})=(1+2t)e^{2t+x}$. The polynomial factor changes, the exponent stays the same.

\medskip\textbf{Step 2: what is $\ring w$?} $\ring w=\spn_\C\{\partial_t^j\partial_x^kw:j,k\ge0\}$, the vector space spanned by \emph{all} partial derivatives of $w$ (including $w$ itself): all finite complex linear combinations of them.
\begin{itemize}[nosep]
\item \emph{Pure exponential.} $w=e^{2t+3x}$: $\partial_t^j\partial_x^kw=2^j3^kw$, so $\ring w=\spn\{w\}$ and $\dim_\C\ring w=1$.
\item \emph{Polynomial times exponential.} $w=te^x$: $\partial_tw=e^x$, $\partial_t^2w=0$, $\partial_xw=te^x$; every higher derivative is a combination of $te^x$ and $e^x$, so $\ring w=\spn\{te^x,e^x\}$ and $\dim_\C\ring w=2$. Although $w$ has infinitely many derivatives, two independent functions span them all.
\end{itemize}

\medskip\textbf{Step 3 ($\Rightarrow$).} Let $f=p\,e^{\tau t+\xi x}$. Then $\partial_tf=(\partial_tp+\tau p)e^{\tau t+\xi x}$ and $\partial_xf=(\partial_xp+\xi p)e^{\tau t+\xi x}$. Differentiating a polynomial never raises its degree, so every derivative of $f$ is $q\,e^{\tau t+\xi x}$ with $\deg q\le\deg p$. The polynomials of degree $\le d$ in two variables form a space of dimension $(d+1)(d+2)/2$, so the derivatives of each term span a finite-dimensional space; $w$ has finitely many terms, so the same holds for $w$.

\medskip\textbf{Step 4 ($\Leftarrow$).} Suppose $W=\ring w$ is finite-dimensional. Differentiation maps derivatives to derivatives, so $\partial_tW\subseteq W$ and $\partial_xW\subseteq W$: they define linear operators $T=\partial_t|_W$ and $X=\partial_x|_W$. They commute ($TX=XT$, mixed partials commute). Commuting complex matrices are simultaneously triangularizable, so in a suitable basis both are upper triangular. Their diagonal entries give pairs of joint eigenvalues $(\tau_j,\xi_j)$; the joint generalized eigenspaces produce functions $p_j(t,x)e^{\tau_jt+\xi_jx}$: the diagonal part of differentiation produces the exponential, the nilpotent parts produce the polynomial factors (they vanish when raised to a high enough power, since the space is finite-dimensional). There are finitely many pairs, hence $w\in\EP$.

\medskip\textbf{Step 5: $\ring w\cong\ring/\Ann(w)$.} Define $\Phi:\ring\to\ring w$, $Q\mapsto Qw$. It is surjective (by definition of $\ring w$) and its kernel is exactly $\Ann(w)$ (since $\Phi(Q)=0$ means $Qw=0$). By the first isomorphism theorem $\ring w\cong\ring/\Ann(w)$. The quotient identifies two operators exactly when they act identically on $w$: $P\sim Q\iff(P-Q)w=0$. Hence
\[
w\in\EP\iff\dim_\C\ring/\Ann(w)<\infty .
\]

\medskip\textbf{Step 6: $\Var(\Ann w)$ is the set of exponents.} For $w=te^x$, $\Ann(w)=(\partial_t^2,\partial_x-1)$, with symbol ideal $(\tau^2,\xi-1)$. Its zero set is $\{\tau^2=0,\ \xi=1\}=\{(0,1)\}$ (over $\C$, $\tau^2=0$ forces $\tau=0$), exactly the exponent pair of $te^{0t+1x}$. For a general exponential-polynomial the variety gives the \emph{distinct} exponent pairs that occur. It does \emph{not} record multiplicities or the full polynomial prefactors: those affect the structure of the annihilator but introduce no new exponent pairs.

\medskip\textbf{Step 7: why it is computationally important.} The criterion converts a question about the \emph{form} of a function into finite-dimensional linear algebra. If one can establish relations such as $\partial_t^2w=2\partial_tw+w$ and $\partial_xw=3w$, then every mixed derivative reduces to a combination of $w$ and $\partial_tw$, so $\dim\ring w\le2$ and $w\in\EP$. The big picture: the annihilator encodes all differential equations satisfied by $w$; the quotient $\ring/\Ann(w)$ counts the independent derivatives; the joint eigenvalues of differentiation identify the exponential modes.

\subsection{The ODE prototype, step by step}\label{expl:E4}
\noindent\emph{(E4: end of Section 3, before the keypoint)}
This section is the one-dimensional version of the whole theory. It starts with an ODE because in one variable the algebra is completely explicit: roots of a polynomial determine the solutions, and the roots lead naturally to the minimal annihilator, the companion matrix and the spectral projectors.

\medskip\textbf{Step 1: the ODE.} $p(D)u=0$ with $D=d/dx$ and $p$ of degree $m$, factored over $\C$ as $p(\lambda)=\prod_{j=1}^r(\lambda-\lambda_j)^{m_j}$, $\sum_jm_j=m$: the $\lambda_j$ are the distinct roots, $m_j$ their multiplicities, $m$ the total degree. For $p(\lambda)=(\lambda-1)^2(\lambda+2)^3$ the equation $(D-1)^2(D+2)^3u=0$ has order five.

\medskip\textbf{Step 2: why solutions are $q(x)e^{\lambda x}$.} First, $(D-\lambda)e^{\lambda x}=0$. If the root has multiplicity greater than one, further solutions appear: $(D-\lambda)(xe^{\lambda x})=e^{\lambda x}$, hence $(D-\lambda)^2(xe^{\lambda x})=0$; in general $(D-\lambda)^m(x^ke^{\lambda x})=0$ for $0\le k<m$. The reason is the identity
\[
(D-\lambda)\big(q(x)e^{\lambda x}\big)=q'(x)e^{\lambda x}:
\]
each application of $D-\lambda$ differentiates the polynomial prefactor once, and after $m$ applications a polynomial of degree $<m$ has vanished. A root of multiplicity $m_j$ therefore contributes the $m_j$ solutions $e^{\lambda_jx},xe^{\lambda_jx},\dots,x^{m_j-1}e^{\lambda_jx}$, and the complete solution space is
\[
\ker p(D)=\Big\{\sum_{j=1}^rq_j(x)e^{\lambda_jx}:\deg q_j<m_j\Big\}.
\]
\emph{Worked example.} For $(\lambda-1)^2(\lambda+2)^3$: root $1$ (multiplicity 2) gives $e^x,xe^x$; root $-2$ (multiplicity 3) gives $e^{-2x},xe^{-2x},x^2e^{-2x}$. Five independent solutions, so the solution space has dimension five $=m$: the number of natural basis functions is finite and equals the degree of $p$.

\medskip\textbf{Step 3: the minimal annihilator.} For a particular function $u$, $\Ann(u)=\{Q\in\C[D]:Q(D)u=0\}$. In one variable the annihilator of a nonzero $u$ is generated by a single polynomial of smallest degree; normalized to be monic it is the \emph{minimal annihilating polynomial} of $u$. Example: $u=xe^x$ has $(D-1)u=e^x$ and $(D-1)^2u=0$, so its minimal annihilator is $p_u(D)=(D-1)^2$; the multiplicity two tells that the solution has a polynomial prefactor of degree one. The polynomial $p$ defining the ODE need not be minimal for a particular solution: $u=e^x$ solves $(D-1)^2u=0$, but its minimal annihilator is only $D-1$.

\medskip\textbf{Step 4: the companion matrix.} Write $p(\lambda)=\lambda^m+a_{m-1}\lambda^{m-1}+\dots+a_0$. Then $u^{(m)}=-a_{m-1}u^{(m-1)}-\dots-a_0u$. With the state vector $Y=(u,u',\dots,u^{(m-1)})^\top$ the ODE becomes $Y'=CY$, with $C$ the companion matrix displayed in the text. Its characteristic polynomial is $p(\lambda)$: the roots of the polynomial become the eigenvalues of a matrix, and the multiplicities determine the possible Jordan blocks.

\medskip\textbf{Step 5: spectral projectors.} Suppose first the roots are distinct. The Lagrange polynomials $\ell_j(z)=\prod_{k\ne j}\frac{z-\lambda_k}{\lambda_j-\lambda_k}$ satisfy $\ell_j(\lambda_i)=1$ if $i=j$ and $0$ if $i\ne j$. Substituting the matrix gives the projectors $\Pi_j=\ell_j(C)$; each projects onto the eigenspace of $\lambda_j$, and $\Pi_i\Pi_j=0$ for $i\ne j$, $\sum_j\Pi_j=\Id$. Consequently $e^{xC}=\sum_je^{\lambda_jx}\Pi_j$ when $C$ is diagonalizable with distinct eigenvalues. When roots are repeated, eigenspace projectors alone do not capture the polynomial factors: one needs the generalized eigenspaces and the nilpotent parts of the Jordan blocks, which produce $xe^{\lambda x}$, $x^2e^{\lambda x}$, \dots

\medskip\textbf{Step 6: why this is the prototype.}
\begin{center}\small
\begin{tabular}{@{}L{5.6cm}L{8.6cm}@{}}
\toprule
One-dimensional ODE & Later multivariable theory\\\midrule
polynomial $p(\lambda)$ & polynomial symbol $P(\tau,\xi)$\\
roots $\lambda_j$ & exponent pairs $(\tau_j,\xi_j)$\\
minimal annihilating polynomial & annihilator ideal and its gcd\\
companion matrix & commuting differentiation operators on a finite-dimensional space\\
Lagrange spectral projectors & projectors isolating the individual exponential modes\\
\bottomrule
\end{tabular}
\end{center}
In one variable the roots are isolated points and polynomial factorization gives the whole picture. In two variables the common zero set of polynomials can contain entire curves, so the geometry and the ideal decomposition become more subtle.

\subsection{What Theorem~\ref{thm:ident} asks, and why}\label{expl:E5}
\noindent\emph{(E5: Section 4.1, after the easy-direction proof / before the takeaway)}
\textbf{Step 1: the problem.} Given an open connected $\Omega\subset\R^2$, a distribution $u\in\mathcal D'(\Omega)$ and a constant-coefficient operator with $P(\partial)u=0$, does $u$ \emph{uniquely identify} $P$, in the sense $\Ann(u)=(P)$? That equality says that every constant-coefficient operator annihilating $u$ is a multiple of $P$: there is no genuinely smaller equation. A single function can satisfy many different PDEs, so the answer is not automatically yes.

\medskip\textbf{Counterexample.} $P=\partial_x^2-1=(\partial_x-1)(\partial_x+1)$ and $u=e^x$: $P(\partial)u=0$, but also $(\partial_x-1)u=0$, so $\Ann(u)\ne(P)$ (in $\C[\partial_x]$ it equals $(\partial_x-1)$; in $\ring$ it is $(\partial_t,\partial_x-1)$). Knowing $u=e^x$ does not tell us that the second-order equation is the minimal one.

\medskip\textbf{Step 2: factor $P$.} $P=\prod_{i=1}^rP_i^{e_i}$ with distinct irreducible $P_i$. Example: $P=(\tau-\xi)^2(\tau+\xi)$ has $P_1=\tau-\xi$ ($e_1=2$) and $P_2=\tau+\xi$ ($e_2=1$). Each factor is a potential piece of the equation that might be unnecessary to explain this particular $u$.

\medskip\textbf{Step 3: the key construction.} For each $i$ put $w_i=(P/P_i)(\partial)u$. Because $P(\partial)u=0$,
\[
P_i(\partial)w_i=P_i(\partial)\tfrac{P}{P_i}(\partial)u=P(\partial)u=0,\qquad\text{i.e.}\qquad w_i\in\ker P_i(\partial).
\]
After removing one factor, the resulting function satisfies the PDE of that factor alone. The theorem asks whether $w_i$ is an exponential-polynomial:
\begin{itemize}[nosep]
\item $w_i\notin\EP$: it has infinitely many independent derivatives, which prevents the factor $P_i$ from being removed from the annihilator;
\item $w_i\in\EP$: it lives in a finite-dimensional derivative space, which signals a possible ambiguity in identifying the equation.
\end{itemize}

\medskip\textbf{Step 4: part (a).} $\Ann(u)=(P)\iff(P/P_i)(\partial)u\notin\EP(\Omega)$ for every $i$. If another annihilator $Q$ were not a multiple of $P$, the annihilator would contain more information than $(P)$; the theorem rules this out precisely by the non-$\EP$ tests. \emph{Check on the counterexample:} for $P_1=\partial_x-1$, $P/P_1=\partial_x+1$ and $w_1=(\partial_x+1)e^x=2e^x\in\EP$; for $P_2=\partial_x+1$, $w_2=(\partial_x-1)e^x=0\in\EP$. The tests say $\EP$, so $u$ does not identify $P$, as it should. Note the theorem concerns the \emph{full annihilator ideal}, not merely whether $P$ is satisfied.

\medskip\textbf{Step 5: part (b).} $\Var(\Ann u)=\Var(P)\iff(P/P_i^{e_i})(\partial)u\notin\EP$ for all $i$. In (a) one removes one copy of $P_i$, in (b) the whole power $P_i^{e_i}$, because the characteristic variety sees zeros, not multiplicities. For instance $(\tau-\xi)^2$ and $\tau-\xi$ have the same zero set $\{\tau=\xi\}$ but are different polynomials. Identifying the variety is therefore weaker than identifying the annihilator.

\medskip\textbf{Step 6: part (c).} If $P$ is \emph{reduced} (all $e_i=1$), removing one copy of $P_i$ equals removing its whole power, so (a) and (b) coincide. If $P$ is monic in $\tau$ of degree $k$ in $\tau$, the test can be done on the finite $t$-jet $u,\partial_tu,\dots,\partial_t^{k-1}u$: a jet is just a finite list of derivatives, so the test needs only finitely many derivative components.

\medskip\textbf{The logical chain.} Given $P(\partial)u=0$ $\to$ factor $P$ $\to$ remove factors, $w_i=(P/P_i)(\partial)u$ $\to$ test membership in $\EP$ (finite-dimensional derivative span, Lemma~\ref{lem:ep}) $\to$ decide whether $u$ determines the full annihilator or only its characteristic variety. One precision: the theorem does \emph{not} say that an exponential-polynomial solution can never identify any PDE; it tests the particular functions $w_i$ obtained after removing factors, not whether $u$ itself is in $\EP$.

\subsection{The point spectrum and the intrinsic operator $G_u$}\label{expl:E6}
\noindent\emph{(E6: Section 4.2, after the aside "Why G\_u and not the full variety")}
\textbf{Step 1: the decomposition.} By Lemma~\ref{lem:gcd}, $\Ann(u)=G_uJ_u$ with $G_u=\Gcd\Ann(u)$ and $\Var(J_u)$ finite. $G_u$ captures the common polynomial factor of all annihilating equations; $J_u$ captures the residual finite structure. The paper calls $\sigma_p(u)=\Var(J_u)$ the \emph{point spectrum} of $u$: a point $(\tau_0,\xi_0)$ corresponds to the exponential $e^{\tau_0t+\xi_0x}$. ``Spectrum'' is used in an algebraic sense: a finite set of exponent pairs, not a set of Fourier frequencies.

\medskip\textbf{Step 2: why $\Ann(u)=(P)\iff\sigma_p(u)=\emptyset$ and $G_u\sim P$.} Here $G_u\sim P$ means they differ by a nonzero constant.
\begin{itemize}[nosep]
\item Suppose $\sigma_p(u)=\emptyset$, i.e.\ $\Var(J_u)=\emptyset$. Over $\C$ the Nullstellensatz says an ideal without common zeros is the whole ring, so $J_u=\ring$ and $\Ann(u)=G_u\ring=(G_u)$: with empty point spectrum the annihilator is generated by its gcd. If moreover $G_u\sim P$, then $\Ann(u)=(P)$.
\item Conversely if $\Ann(u)=(P)$, the gcd is $P$ up to a constant, the cofactor ideal is $\ring$ and its zero set is empty.
\end{itemize}
Interpretation: the equation is fully identifiable when there is no residual finite point spectrum and the common gcd agrees with the proposed operator.

\medskip\textbf{Step 3: a single exponential.} $u=e^{\tau_0t+\xi_0x}$ has $\partial_tu=\tau_0u$, $\partial_xu=\xi_0u$, so $\Ann(u)=(\tau-\tau_0,\xi-\xi_0)$. The generators have no common nonconstant factor: $G_u=1$, $J_u=(\tau-\tau_0,\xi-\xi_0)$, and
\[
\sigma_p(u)=\{(\tau_0,\xi_0)\}:
\]
one point, exactly the exponential in $u$. Since $G_u=1$ and the spectrum is nonempty, this $u$ does not identify any nonconstant operator as the generator of its full annihilator: $\partial_t-\tau_0$ and $\partial_x-\xi_0$ both annihilate $u$, and neither alone generates $\Ann(u)$.

\medskip\textbf{Step 4: the meaning of $\ring u\cap\EP\cong\ring/J_u$.} $\ring u$ is the span of all derivatives of $u$; $\ring u\cap\EP$ consists of the exponential-polynomials obtainable from $u$ by constant-coefficient operators and linear combinations; the quotient $\ring/J_u$ describes this part algebraically. For the single exponential, $\ring/J_u\cong\C$, matching $\ring u=\spn\{u\}$, one-dimensional and entirely inside $\EP$. In general the finite-dimensional quotient also records \emph{multiplicity} information that the set $\sigma_p(u)$ alone loses.

\medskip\textbf{Step 5: why $G_u$ is the intrinsic operator.} Knowing some $P$ with $P(\partial)u=0$, one wants the operator intrinsic to $u$ rather than $P$ itself, which may contain unnecessary factors. Strategy: (1) factor $P=\prod P_i^{e_i}$; (2) remove factors and test whether the resulting functions are in $\EP$; (3) recover $G_u$ and $J_u$; (4) if $J_u=\ring$, the point spectrum is empty and the annihilator is generated by $G_u$. At most $\sum_ie_i$ membership tests are needed.

\medskip\textbf{Step 6: the repeated-factor example.} The variety $\Var(\Ann u)$ can replace $G_u$ only for reduced $P$, because zero sets do not see multiplicities. The aside above uses $P=\tau^2$ and $u=e^{x^2}$: here $\Ann(u)=(\tau)$, so $G_u=\tau\not\sim P$ while $\Var(\Ann u)=\Var(P)=\{\tau=0\}$. The variety criterion would wrongly say that $u$ identifies $\tau^2$. (An earlier draft used $P=\tau$ for this example, which is reduced and so shows nothing about repeated factors; the revised text uses $\tau^2$.)

\subsection{The d'Alembert function does \emph{not} split by terms (Proposition~\ref{prop:terms})}\label{expl:E7a}
\noindent\emph{(E7a: Section 5.1 (after the heat example, before 5.2))}
\textbf{Enunciation.} If $L=L_1+L_2$ (constant coefficients), then
\[
\ker L\cap(\ker L_1+\ker L_2)\ \subseteq\ \ker(L,\ L_1L_2),
\]
and this space is finite-dimensional as soon as $\Var(L_1)\cap\Var(L_2)$ is finite, because $\Var(L,L_1L_2)=\Var(L_1)\cap\Var(L_2)$. In words: only a finite-dimensional family of solutions of $L$ can be written as (solution of the first term) $+$ (solution of the second term).

\medskip\textbf{Applying it to the wave equation.} Write the wave operator as a sum of its two \emph{terms}, $L=\partial_t^2-c^2\partial_x^2=L_1+L_2$ with $L_1=\partial_t^2$ and $L_2=-c^2\partial_x^2$. Then
\[
\ker L_1=\{a(x)+t\,b(x)\},\qquad \ker L_2=\{\alpha(t)+x\,\beta(t)\},
\]
so every element of $\ker L_1+\ker L_2$ is killed by $L_1L_2=-c^2\partial_t^2\partial_x^2$ (the operators commute: $L_1L_2u=L_2L_1a+L_1L_2b=0$).

\medskip\textbf{Step 1: a d'Alembert solution.} Take $u=f(x+ct)+g(x-ct)$, which solves $Lu=0$. Computing,
\[
\partial_t^2\partial_x^2f(x+ct)=c^2f''''(x+ct),\qquad \partial_t^2\partial_x^2g(x-ct)=c^2g''''(x-ct),
\]
so $L_1L_2u=-c^4\big(f''''(x+ct)+g''''(x-ct)\big)$.

\medskip\textbf{Step 2: when is it zero?} $f''''(s)+g''''(r)=0$ for all independent $s=x+ct$, $r=x-ct$ forces $f''''\equiv k$ and $g''''\equiv-k$ for a constant $k$. So $f$ and $g$ must be polynomials (of degree $\le4$, with matching quartic coefficients). The d'Alembert solutions that lie in $\ker L_1+\ker L_2$ are polynomials only; a generic one, such as $f(s)=e^s$ ($u=e^{x+ct}$ for $g=0$), gives $L_1L_2u=-c^4u\ne0$, so it is \emph{not} a sum of a solution of $\partial_t^2u=0$ and one of $\partial_x^2u=0$.

\medskip\textbf{Step 3: the finite dimension.} Here $\Var(L_1)\cap\Var(L_2)=\{\tau^2=0\}\cap\{\xi^2=0\}=\{(0,0)\}$, a point, so the space $\ker(L,L_1L_2)$ is finite-dimensional. It is the polynomial family above: $f,g$ cubic polynomials plus the shared quartic term $\pm k\,s^4/24$, modulo the constant overlap $(f+\kappa,g-\kappa)$; this gives $4+4-1+1=8$ parameters, which agrees with the intersection multiplicity $\dim\ring/(\tau^2-c^2\xi^2,\tau^2\xi^2)=4+4=8$.

\medskip\textbf{Conclusion.} The wave solution $u=f(x+ct)+g(x-ct)$ is the sum of \emph{factor} components (next subsection) but, in general, not of \emph{term} components. This is the structural reason why any attempt to split by terms leaves a remainder (Section~\ref{sec:inside}).

\subsection{The d'Alembert function proves Theorem~\ref{thm:factor}, step by step}\label{expl:E7b}
\noindent\emph{(E7b: Section 5.2, after the Takeaway of Theorem 5.3)}
\textbf{Enunciation (Factor decomposition).} Let $\Omega\subset\R^2$ be bounded, connected and simply connected, and $P=P_1^{m_1}\cdots P_N^{m_N}$ with $Q_j=P_j^{m_j}$ pairwise coprime. Then every $u\in\ker P$ can be written $u=\sum_{j=1}^Nu_j$ with $u_j\in\ker Q_j(\partial)$, and the ambiguity of the decomposition is an exponential-polynomial space of dimension exactly $\delta(P)=\sum_{i<j}\dim\ring/(Q_i,Q_j)$.

\medskip\textbf{The case.} Wave equation, $N=2$, $m_1=m_2=1$:
\[
P=\tau^2-c^2\xi^2=Q_1Q_2,\qquad Q_1=\partial_t-c\partial_x,\quad Q_2=\partial_t+c\partial_x,
\]
and the d'Alembert solution $u=f(x+ct)+g(x-ct)$. We check each claim of the theorem.

\medskip\textbf{Step 1: the hypotheses.} $Q_1,Q_2$ are distinct linear polynomials, hence coprime (no common factor). $\Omega$ is any bounded simply connected region, e.g.\ a disc.

\medskip\textbf{Step 2: the kernels.} $Q_1f(x+ct)=cf'-cf'=0$ and $Q_2g(x-ct)=-cg'+cg'=0$, so $\ker Q_1=\{f(x+ct)\}$ (a profile moving left) and $\ker Q_2=\{g(x-ct)\}$ (a profile moving right).

\medskip\textbf{Step 3: existence of the splitting.} Take any $u\in\ker P$. Then $w_2=Q_2u\in\ker Q_1$ and $w_1=Q_1u\in\ker Q_2$ (because $Q_1Q_2u=0$). Since $Q_2-Q_1=2c\,\partial_x$,
\[
\partial_xu=\tfrac1{2c}\big(Q_2u-Q_1u\big).
\]
Integrating in $x$ at fixed $t$, define $u_1=\tfrac1{2c}\int^xQ_2u\,\dd x'$ and $u_2=-\tfrac1{2c}\int^xQ_1u\,\dd x'$. If $Q_2u=h(x+ct)$, then $u_1=\tfrac1{2c}H(x+ct)$ with $H'=h$, and $Q_1u_1=0$; likewise $u_2\in\ker Q_2$. By construction $\partial_x(u-u_1-u_2)=0$, so $u-u_1-u_2=k(t)$. Applying $P$ gives $k''(t)=0$, so $k(t)=a+bt$; and $t=\tfrac1{2c}\big((x+ct)-(x-ct)\big)$ is itself a sum of an element of $\ker Q_1$ and one of $\ker Q_2$. Hence $u=u_1+u_2$ with $u_j\in\ker Q_j$ after absorbing $k$: the splitting exists.

\medskip\textbf{Step 4: the explicit example.} With $c=1$, $u=\dfrac1{1+(x+t)^2}+e^{-(x-t)^2}$: $Q_2u=2f'(x+t)$ and $Q_1u=-2g'(x-t)$, so $u_1=\tfrac12\int^x2f'(x'+t)\dd x'=f(x+t)$ and $u_2=-\tfrac12\int^x(-2g')\dd x'=g(x-t)$, with $k=0$: the recovered components are exactly the two moving profiles.

\medskip\textbf{Step 5: the ambiguity has dimension $\delta(P)=1$.} $Q_2-Q_1=2c\,\xi$ and $Q_1+Q_2=2\tau$, so $(Q_1,Q_2)=(\tau,\xi)$, which is not the unit ideal (the two lines meet at the origin), and $\ring/(\tau,\xi)\cong\C$. So $\delta(P)=\dim\ring/(Q_1,Q_2)=1$. Directly: $f(x+ct)+g(x-ct)=0$ for all $(t,x)$ forces $f\equiv\kappa$, $g\equiv-\kappa$ for a constant $\kappa$. The only freedom is $(f,g)\mapsto(f+\kappa,g-\kappa)$: a one-dimensional space of constants, which are exponential-polynomials -- exactly the dimension the theorem predicts.

\medskip\textbf{Step 6: where the theorem's tools appear.} The finite quotient $\ring/(Q_1,Q_2)$ measures the failure of the one-variable picture. In one variable coprime polynomials satisfy a B\'ezout identity $aQ_1+bQ_2=1$ and $u=bQ_2u+aQ_1u$ splits explicitly with a direct sum; in two variables only $(Q_1,Q_2)\supseteq(\tau,\xi)$ holds, and that is what produces the constant ambiguity. The wave case is the simplest example where the proof by exact sequence $0\to\ring/(Q_1Q_2)\to\ring/(Q_1)\oplus\ring/(Q_2)\to\ring/(Q_1,Q_2)\to0$ can be followed by hand.

\medskip\textbf{Contrast with 5.1.} The same function is a sum of \emph{factor} components for every datum $(f,g)$, but a sum of \emph{term} components only when $f,g$ are polynomials. This is the difference between a factor and a term.

\subsection{The Lagrange projector for evolution factors}\label{expl:E8}
\noindent\emph{(E8: Section 5.3, after Eq. (lagrange))}
\textbf{Step 1: the evolution operator.} $P(\tau,\xi)=\prod_{j=1}^N(\tau-a_j(\xi))$ with $a_j$ distinct polynomials in $\xi$. Replacing $\tau,\xi$ by $\partial_t,\partial_x$ the equation is $\prod_j(\partial_t-a_j(\partial_x))u=0$. Example: $a_1(\xi)=\xi$, $a_2(\xi)=-\xi$ gives $(\partial_t-\partial_x)(\partial_t+\partial_x)u=0$, the one-dimensional wave equation $u_{tt}-u_{xx}=0$. Because the factors are distinct, a solution decomposes as $u=\sum_ju_j$ where each component satisfies its own first-order law $\partial_tu_j=a_j(\partial_x)u_j$, i.e.\ $(\partial_t-a_j(\partial_x))u_j=0$.

\medskip\textbf{Step 2: isolating a component.} Fix $j$ and apply $Q_j=\prod_{k\ne j}(\partial_t-a_k(\partial_x))$ to $u=\sum_lu_l$. Constant-coefficient operators commute, so $Q_ju=\sum_lQ_ju_l$.
\begin{itemize}[nosep]
\item \emph{Case $l\ne j$.} $Q_j$ contains the factor $\partial_t-a_l(\partial_x)$, which kills $u_l$: $Q_ju_l=0$. Every component except $u_j$ disappears.
\item \emph{Case $l=j$.} No factor of $Q_j$ kills $u_j$; use $\partial_tu_j=a_j(\partial_x)u_j$: for each $k\ne j$, $(\partial_t-a_k(\partial_x))u_j=(a_j-a_k)(\partial_x)u_j$. Applying all factors, $Q_ju=\prod_{k\ne j}(a_j-a_k)(\partial_x)\,u_j$.
\end{itemize}
Hence
\[
u_j=\Big[\prod_{k\ne j}(a_j-a_k)(\partial_x)\Big]^{-1}\prod_{k\ne j}(\partial_t-a_k(\partial_x))\,u,
\]
which is Eq.~\eqref{eq:lagrange}, provided the inverse operator exists on the function space considered. It is not automatically a differential operator: it is in general a Fourier multiplier, possibly singular at frequencies where $a_j=a_k$.

\medskip\textbf{Step 3: why this is Lagrange interpolation.} The Lagrange polynomials $L_j(z)=\prod_{k\ne j}\frac{z-a_k}{a_j-a_k}$ satisfy $L_j(a_l)=1$ for $l=j$, $0$ otherwise; for a matrix $C$ diagonalized in eigenvectors, $L_j(C)$ acts as the identity on the $j$-th eigenspace and as zero on the others. Here the role of $C$ is played by $\partial_t$ and the eigenvalues $a_j$ are replaced by the operators $a_j(\partial_x)$:
\[
\Pi_j=\prod_{k\ne j}\frac{\partial_t-a_k(\partial_x)}{a_j(\partial_x)-a_k(\partial_x)},\qquad \Pi_ju=u_j .
\]
The analogy is exact at the algebraic level: $\Pi_j$ is the projector that extracts the $j$-th evolution component.

\medskip\textbf{Step 4: rational in $\xi$.} With the Fourier convention $\partial_x\mapsto i\kappa$ ($\kappa$ real frequency), $a_j(\partial_x)$ becomes multiplication by $a_j(i\kappa)$ and the multiplier is
\[
\ell_j(\kappa,\partial_t)=\prod_{k\ne j}\frac{\partial_t-a_k(i\kappa)}{a_j(i\kappa)-a_k(i\kappa)} .
\]
The denominators vanish exactly where $a_j(i\kappa)=a_k(i\kappa)$: these are the potential poles of the projector. Strictly, the rationality concerns the dependence on the spatial frequency; the projector still involves $\partial_t$ before it is restricted to a solution component, on which $\partial_t$ acts as multiplication by $a_j(i\kappa)$.

\medskip\textbf{Step 5: the wave equation.} $(\partial_t-\partial_x)(\partial_t+\partial_x)u=0$, $u=u_1+u_2$ with $\partial_tu_1=\partial_xu_1$ and $\partial_tu_2=-\partial_xu_2$. Applying $\partial_t+\partial_x$: the $u_2$ term vanishes and $(\partial_t+\partial_x)u_1=2\partial_xu_1$, so
\[
u_1=\tfrac12\,\partial_x^{-1}(\partial_t+\partial_x)u,\qquad u_2=-\tfrac12\,\partial_x^{-1}(\partial_t-\partial_x)u .
\]
In Fourier variables $\partial_x^{-1}$ is multiplication by $1/(i\kappa)$: a pole at $\kappa=0$. This does not mean the wave equation has no decomposition (Section~\ref{sec:groups} proves it does), only that this particular formula needs a choice of inverse (an integration constant) or extra assumptions at zero frequency.

\medskip\textbf{Step 6: what ``coprime'' means.} $\tau-a_j(\xi)$ and $\tau-a_k(\xi)$ are coprime in $\C[\tau,\xi]$ exactly when $a_j-a_k$ is not the zero polynomial; for distinct $a_j$ this holds, though $a_j-a_k$ may still vanish at particular frequencies, creating poles. Three levels must be kept apart: \emph{coprime factors} (the algebraic decomposition exists in the localized, rational-operator sense); \emph{no poles at any frequency} (stronger, not guaranteed by coprimality); \emph{a bounded or well-defined projector on a chosen function space} (needs analytic conditions near frequencies where branches meet). The message: distinct evolution branches admit algebraic projectors, like distinct eigenvalues of a matrix; in spatial Fourier variables the projectors are rational multipliers whose denominators reveal where the branches intersect.

\subsection{Regularity, the sum map $S$, $\gamma_N$ and the Friedrichs angle}\label{expl:E9}
\noindent\emph{(E9: Section 5.5, just before the final Takeaway of Section 5)}
\textbf{The guiding idea.} $\gamma_N$ measures \emph{how reliably a solution can be separated into components belonging to different kernel spaces, and how much errors can be amplified when this is done}. Being able to decompose algebraically does not guarantee that the pieces have the regularity of the original solution; the obstruction is geometric: the component spaces may be too close to one another, which makes the decomposition unstable.

\medskip\textbf{Step 1: why a component can lose a derivative.} $Pu=0$, $u=\sum_ju_j$, $Q_ju_j=0$. The pieces exist, but $u\in H^m$ does not imply $u_j\in H^m$, because the formulas that extract $u_j$ may involve inverse differential operators, which can lose regularity. In the wave example $u_1=\tfrac12\partial_x^{-1}(\partial_t+\partial_x)u$: even if $u$ has $m$ Sobolev derivatives the formula alone does not guarantee that $u_1$ has $m$, because of the inverse and its behaviour at problematic frequencies. This is the content of Proposition~\ref{prop:loss}.

\medskip\textbf{Step 2: the component spaces and the sum operator.} $M_j=\ker_{H^m}Q_j$ (admissible components); for two components $M=\ker_{H^m}Q_1$, $N=\ker_{H^m}Q_2$. The decomposition question is packaged in
\[
S:M_1\times\cdots\times M_N\to H^m(\Omega),\qquad S(u_1,\dots,u_N)=\sum_ju_j,\qquad \|(u_1,\dots,u_N)\|^2=\sum_j\|u_j\|_{H^m}^2 .
\]
Recovering the components means inverting $S$ on the relevant solution space. Two things can go wrong: \emph{non-uniqueness} (different tuples with the same sum: $\ker S=\{(u_j):\sum u_j=0\}$; if $v\in M\cap N$ then $(v,-v)\in\ker S$, so overlapping spaces give non-uniqueness) and \emph{instability} (even if unique, a tiny change in the sum may require a large change in the pieces).

\medskip\textbf{Step 3: what $\gamma_N$ is.}
\[
\gamma_N=\inf\sigma\big(S^*S|_{(\ker S)^\perp}\big)=\inf_{\substack{v\in(\ker S)^\perp\\v\ne0}}\frac{\|Sv\|_{H^m}^2}{\|v\|^2}.
\]
$S^*$ is the adjoint, $S^*S$ is positive self-adjoint, and $(\ker S)^\perp$ removes the directions that change the pieces without changing the sum (without this removal the ratio would be $0$ regardless of the stability of the genuinely recoverable part). If $\gamma_N>0$ then $\|Sv\|^2\ge\gamma_N\|v\|^2$ on $(\ker S)^\perp$, i.e.
\[
\|v\|\le\gamma_N^{-1/2}\|Sv\|:
\]
this is the stability estimate; the smaller $\gamma_N$, the larger the possible amplification. \emph{Technical caveat:} if the restriction has continuous spectrum accumulating at $0$, ``smallest nonzero eigenvalue'' (used loosely in the text after Figure~\ref{fig:friedrichs}) is not literally correct; the robust quantity is the \emph{infimum of the spectrum} on $(\ker S)^\perp$, as in the theorem.

\medskip\textbf{Step 4: two components and the Friedrichs angle.} Let $I=M\cap N$. The Friedrichs cosine is $c_F=\sup\{|\langle x,y\rangle|:x\in M\cap I^\perp,\ y\in N\cap I^\perp,\ \|x\|=\|y\|=1\}$, and $\theta\in[0,\pi/2]$ is defined by $c_F=\cos\theta$. If $c_F$ is close to $1$ the spaces have almost parallel directions; if it is clearly below $1$ they are well separated; $c_F<1$ means the sum $M+N$ is closed. A finite-dimensional analogy: decomposing a vector along two nearly parallel directions gives large components that nearly cancel, leaving a small sum; recovering them from the small sum is numerically unstable. The paper proves $\gamma_2=1-c_F$, so $c_F<1\iff\gamma_2>0$.

\medskip\textbf{Step 5: why pairwise angles do not suffice for $N\ge3$.} Three vectors in a plane pointing in three different directions can satisfy $v_1+v_2+v_3=0$: the instability depends on combinations of all components, not only on pairs. $S$ captures these collective effects; its smallest nonzero singular value is $\sqrt{\gamma_N}$, so $1/\sqrt{\gamma_N}$ is the natural condition-number-like factor for recovering components from their sum, as in a least-squares problem where a small singular value makes inversion sensitive to errors.

\medskip\textbf{Step 6: density.} The theorem needs $M_1+\dots+M_N$ to be both closed and dense in $\ker_{H^m}P$. Dense plus closed gives $M_1+\dots+M_N=\ker_{H^m}P$, i.e.\ a decomposition for every solution, with stability from the positive-angle condition. The text states density on bounded convex domains for every $P$ and $m$ (H\"ormander, Stein; see its status label), so closedness is the decisive remaining condition there.

\medskip\textbf{The chain in six lines.} (1) Algebra produces components attached to the factors. (2) Sobolev asks whether they remain in $H^m$. (3) This is controlled by $S$. (4) $\gamma_N>0$ makes $S$ bounded below after removing its kernel: stable recovery. (5) For $N=2$ this is $c_F<1$. (6) For $N\ge3$ the collective $\gamma_N$, not the pairwise angles, controls the conditioning. One distinction: $\gamma_N>0$ is about the stability of the decomposition map; it is not an algebraic statement about factorization. The factorization says which spaces to consider, their geometry in $H^m$ says whether the components can be recovered with the desired regularity.

\subsection{The additive interaction remainder and why the terms commute}\label{expl:E11}
\noindent\emph{(E11: Section 6.2, after the proof of Theorem 6.1 (before the concrete case))}
\textbf{Step 1: the setup.} $\partial_tu=Au$, $u(0)=u_0$, with $A=\sum_{n=1}^NA_n$. The solution is formally $u(t)=e^{tA}u_0$. If term $n$ acted alone it would produce $S_n(t)=e^{tA_n}$, with change $\Delta_nu_0=(S_n(t)-\Id)u_0$ ($\Delta_n$ is the \emph{increment operator}). The additive interaction remainder is
\[
\Eint^{\mathrm{add}}(t)=u(t)-u_0-\sum_{n=1}^N\Delta_nu_0,
\]
that is, \emph{interaction $=$ full evolution $-$ initial state $-$ sum of the individual increments.}
Note that this is an additive decomposition, \emph{not} a decomposition into kernel spaces as in Section~\ref{sec:groups}.

\medskip\textbf{Step 2: why commutation gives an exact formula.} If $A_iA_j=A_jA_i$ then (for suitable operators and a well-defined evolution) $e^{t\sum_nA_n}=\prod_ne^{tA_n}=\prod_nS_n(t)$. Since $S_n=\Id+\Delta_n$,
\[
u(t)=\prod_{n=1}^N(\Id+\Delta_n)u_0,\qquad
\Eint^{\mathrm{add}}(t)=\Big[\prod_{n=1}^N(\Id+\Delta_n)-\Id-\sum_{n=1}^N\Delta_n\Big]u_0 .
\]
This is pure algebra: multiply the increments, then subtract the constant term and all terms involving a single increment.

\medskip\textbf{Step 3: $N=2$.} $u(t)=S_1S_2u_0=(\Id+\Delta_1)(\Id+\Delta_2)u_0=(\Id+\Delta_1+\Delta_2+\Delta_1\Delta_2)u_0$. Subtracting $u_0$ and the two single increments leaves
\[
\Eint^{\mathrm{add}}(t)=\Delta_1\Delta_2u_0=(S_1-\Id)(S_2-\Id)u_0 .
\]
It is an interaction term because it contains the effect of both increments acting together; if either increment is zero on the relevant state, the interaction vanishes.

\medskip\textbf{Step 4: $N=3$ shows the general structure.}
\[
(\Id+\Delta_1)(\Id+\Delta_2)(\Id+\Delta_3)=\Id+\Delta_1+\Delta_2+\Delta_3+\Delta_1\Delta_2+\Delta_1\Delta_3+\Delta_2\Delta_3+\Delta_1\Delta_2\Delta_3,
\]
so $\Eint^{\mathrm{add}}=(\Delta_1\Delta_2+\Delta_1\Delta_3+\Delta_2\Delta_3+\Delta_1\Delta_2\Delta_3)u_0$: three \emph{pairwise} interactions and one \emph{three-way} interaction. In general $\Eint^{\mathrm{add}}=\sum_{J\subseteq\{1,\dots,N\},|J|\ge2}\big(\prod_{j\in J}\Delta_j\big)u_0$: the empty subset gives the initial state, subsets of size one the individual increments, everything left involves at least two operators.

\medskip\textbf{Step 5: the Fourier formula.} If each $A_n$ is a constant-coefficient operator it becomes multiplication by its symbol $a_n(\kappa)$, so $\widehat{S_n(t)u_0}=e^{ta_n(\kappa)}\hat u_0$ and $\widehat{\Delta_nu_0}=(e^{ta_n(\kappa)}-1)\hat u_0$. Then
\[
\widehat{\Eint^{\mathrm{add}}}(t,\kappa)=\Big[\prod_ne^{ta_n(\kappa)}-1-\sum_n\big(e^{ta_n(\kappa)}-1\big)\Big]\hat u_0(\kappa).
\]
At each frequency this is ordinary scalar algebra; the operators commute because their Fourier symbols are multiplications. \emph{Notation note:} here $\kappa$ is the real frequency and $a_n(\kappa)$ the symbol of $A_n$ in the variable $i\kappa$ (e.g.\ $a(\kappa)=-ic\kappa-D\kappa^2$ for advection--diffusion).

\medskip\textbf{Step 6: small-time reading.} $\Delta_n=e^{tA_n}-\Id=tA_n+O(t^2)$. For two operators, $\Eint^{\mathrm{add}}=\Delta_1\Delta_2u_0=t^2A_1A_2u_0+O(t^3)$ (given enough regularity). Each single increment starts at order $t$, their interaction at order $t^2$: it needs both effects to contribute. For $N$ operators the interaction involving all $N$ increments starts at $t^N$, while the full remainder also contains pairwise interactions, generally starting at $t^2$.

\medskip\textbf{Step 7: why the terms commute.} Two operators that are polynomials in the \emph{same} $\partial_x$, e.g.\ $A_1=a_0+a_1\partial_x+a_2\partial_x^2$ and $A_2=b_0+b_1\partial_x+b_2\partial_x^2$, commute because powers of $\partial_x$ commute: $[A_1,A_2]=0$. More generally any two constant-coefficient operators, polynomials in $\partial_{x_1},\dots,\partial_{x_d}$, commute; for instance $A_1=\partial_x^2+\partial_y$ and $A_2=3\partial_x\partial_y-\partial_y^2$, because partial derivatives commute on sufficiently smooth functions (and on distributions). So for constant coefficients commutation is not a restrictive hypothesis: it comes almost automatically from the structure of the operators.

\medskip\textbf{The conceptual contrast.} Section~\ref{sec:groups} asked whether a solution can be stably separated into pieces belonging to different kernels. Here the question is different: if the generator is a sum of operators, how much of the resulting evolution cannot be explained by adding their individual changes? When the operators commute the answer is exact: expand the product of the evolution operators, remove the initial state and the single-operator contributions, and what remains is exactly the mixed interaction terms.

\subsection{Burgers: why the shock forms and what happens to the interaction}\label{expl:E13}
\noindent\emph{(E13: Section 7.2, after the paragraph on the regimes / table)}
\textbf{Step 1: why inviscid Burgers forms a shock.} $u_t+uu_x=0$ is nonlinear transport: the velocity at each point is $u$ itself. Along a characteristic starting at $x=\xi$, $dx/dt=u$ and $du/dt=0$, so $u$ stays constant, $u(x(t),t)=u_0(\xi)$. For $u_0=\sin\xi$:
\[
x=\xi+t\sin\xi,\qquad u=\sin\xi .
\]
Differentiating, $\partial x/\partial\xi=1+t\cos\xi$; as long as it is positive everywhere the map $\xi\mapsto x$ is invertible and the solution is smooth. It first vanishes when $1+t\cos\xi=0$; since $\min\cos=-1$, the breaking time is $t^*=1$ at $\xi^*=\pi$. The gradient $u_x=\dfrac{\cos\xi}{1+t\cos\xi}$ has numerator $\to-1$ and denominator $\to0$ at $(t,\xi)=(1,\pi)$, so it blows up. \emph{Before $t=1$} characteristics do not cross and there is one value of $u$ per $x$. \emph{At and after} $t=1$ they compress and cross, the parametric curve folds, no single-valued classical solution exists, and the entropy solution replaces the fold with a shock (for this symmetric datum at $x=\pi$).

\medskip\textbf{Step 2: where the interaction enters.} Viscous Burgers $u_t+uu_x=\nu u_{xx}$ is split into nonlinear transport $A_{\mathrm{tr}}u=-uu_x$ and diffusion $A_{\mathrm{diff}}u=\nu u_{xx}$. With $S(t)$ the full flow and $S_1,S_2$ the single-term flows,
\[
\Eint(t)=S(t)u_0-S_1(t)u_0-S_2(t)u_0+u_0 .
\]
Before the shock all flows can be compared with smooth expansions; after it, the transport flow develops a discontinuity while the viscous solution smooths the transition over a thin layer. That difference drives everything below.

\medskip\textbf{Step 3: why the smooth expansion breaks down.} Before breaking there is a first-order expansion in $\nu$ with remainder $O(\nu^2\varepsilon^2t^3)$ (the exact coefficient depends on the definitions of the paper; the parameter $\varepsilon$ is not defined in the text and is flagged in the author's notes). A perturbative expansion is useful while the relevant solutions stay smooth. At $t^*$ the inviscid gradient is unbounded and estimates relying on bounded derivatives fail. \emph{Precision:} what blows up at $t^*$ is the perturbative expression or its estimates, not the viscous solution itself, which stays finite because positive viscosity regularizes.

\medskip\textbf{Step 4: a jump that survives $\nu\to0$.} After the shock the transport flow is discontinuous while the viscous flow has a thin smooth layer. Even if the viscous profile converges to the entropy solution, the individual pieces of the difference need not converge so that their pointwise jumps cancel; hence $\Eint$ can have a jump of size independent of $\nu$. This does not contradict smallness in $L^1$: a function of amplitude $O(1)$ supported on an interval of width $O(\nu)$ has $\|E\|_{L^\infty}=O(1)$ but $\|E\|_{L^1}=O(\nu)$. The integral measures the area under $|E|$, not the maximum height.

\medskip\textbf{Step 5: the three estimates.}
\begin{center}\small
\begin{tabular}{@{}L{4.3cm}L{3.0cm}L{6.6cm}@{}}
\toprule
Quantity & Claimed scale & Interpretation\\\midrule
$\|\Eint(t)\|_{L^1}$ at fixed $t$ & $\Theta(\nu)$ & small in total absolute magnitude at a fixed post-shock time\\
$\sup_{t\le T}\|\Eint(t)\|_{L^1}$ & $\nu|\log\nu|+O(\nu)$ & the worst case over the time interval has a logarithmic enhancement\\
$\|S-S_1S_2\|_{L^1}$ & $\Theta(\sqrt\nu)$ & a different, larger scale: the error of replacing the full evolution by a composition of separate flows\\
\bottomrule
\end{tabular}
\end{center}
$\Theta(\nu)$ means bounded above and below by positive multiples of $\nu$ as $\nu\to0$. The third quantity is a different object: it concerns $S-S_1S_2$, not the additive remainder, and must not be conflated with it. Exact constants and rates depend on the definitions of the split flows, the domain and the theorem's hypotheses; the status labels in the text say what was proved and what is conditional.

\medskip\textbf{Step 6: the broader lesson.} Three regimes: before the shock, smoothness permits perturbative expansions and derivative-based estimates; at breaking, the inviscid gradient blows up and smoothness-based estimates fail; after the shock, pointwise behaviour can be discontinuous, so $L^1$ estimates become more informative than pointwise bounds. A nonlinear interaction can remain small in an integral norm even when its pointwise behaviour changes dramatically at shock formation.

\subsection{Identifying a nonlinear equation relative to a dictionary}\label{expl:E14}
\noindent\emph{(E14: Section 7.3, after the dictionary criterion)}
\textbf{Step 1: why $\Ann(u)$ is no longer enough.} For a linear equation, e.g.\ $u=e^{x-t}$, $(\partial_t+\partial_x)u=0$ is the action of a linear constant-coefficient operator. In $u_t+uu_x=0$ the term $uu_x$ multiplies two quantities depending on $u$: it is not the action of a linear constant-coefficient operator on $u$, so we need another language: polynomial expressions in the jet variables.

\medskip\textbf{Step 2: jets and the dictionary.} Fix a maximal order $k$ and take the jet variables $u,u_t,u_x,u_{tt},u_{tx},u_{xx},\dots$ up to order $k$. A polynomial in them is a candidate differential relation; Burgers is $F(u,u_t,u_x)=u_t+uu_x$. A \emph{dictionary} $\mathcal D$ is a finite list of candidate monomials, e.g.\ $\{u_t,u_x,u^2,uu_x,u_{xx},\dots\}$, and $V=\spn\mathcal D$ contains all linear combinations. A candidate equation is $c_1u_t+c_2uu_x+c_3u_{xx}=0$; identifying the equation from data means finding the coefficients for which the relation holds on the observed solution. $V$ is a finite-dimensional vector space of candidates, \emph{not} necessarily an ideal: the dictionary limits what one is willing to consider.

\medskip\textbf{Step 3: $I(u)$ versus $[F]$.} $I(u)$ is the differential ideal of \emph{all} polynomial differential relations satisfied by $u$. If $F=u_t+uu_x=0$ holds identically, so do its derivatives, e.g.\ $D_xF=u_{tx}+u_x^2+uu_{xx}$ (total derivative $D_x$). $[F]$, the differential ideal generated by $F$, contains $F$, its differential consequences $D_xF,D_tF,\dots$ and polynomial combinations of them. If $u$ solves $F=0$ then $[F]\subseteq I(u)$, but the inclusion can be strict: a particular solution may satisfy additional relations that do not follow from the equation alone. That is the core of identifiability.

\medskip\textbf{Step 4: why ``$I(u)\cap V=\langle F\rangle$'' is wrong.} If $\mathcal D$ contains the monomials of both $F$ and $D_xF$, both lie in $V$. If $\langle F\rangle$ meant only scalar multiples of $F$, it would not contain $D_xF$, which is certainly a relation satisfied by the solution. The correct comparison is
\[
I(u)\cap V=[F]\cap V:
\]
every relation of the dictionary that the data satisfy is already a differential consequence of the proposed equation. It does \emph{not} require the data to satisfy only the original equation and nothing derived from it.

\medskip\textbf{Step 5: the exact sequence.}
\[
0\to[F]\cap V\to I(u)\cap V\to\mathcal J_u\cap\psi(V)\to0 .
\]
$\mathcal J_u$ is the jet-level obstruction (distinct from the cofactor $J_u$ of Section~\ref{sec:theory}) and $\psi$ the map from candidate expressions to jet-level objects (its precise definition is flagged in the author's notes). The first term holds the dictionary relations explained by the proposed equation, the middle term every dictionary relation the data satisfy, the last the extra relations left after quotienting out the consequences of $F$. Exactness gives $I(u)\cap V=[F]\cap V\iff\mathcal J_u\cap\psi(V)=\{0\}$: the dictionary contains no detectable obstruction. It is a \emph{finite-dictionary} criterion: it does not say the data identify the equation among all conceivable nonlinear PDEs, only relative to the chosen candidate space.

\medskip\textbf{Step 6: shocks and solitons.} Identifiability depends on the solution, not only on the equation. \emph{Burgers with $N$ shocks:} in the paper's evolutionary dictionary the equation is identified iff the maximal polynomial degree satisfies $d\le N$; shock structure gives information that distinguishes the equation from competing polynomial relations, depending on the number of shocks and on the allowed degree of candidate terms (specific to the dictionary and assumptions of the cited result). \emph{$N$-soliton KdV:} the identifying polynomial has total degree $(N+2)2^{N-1}$; for $N=2$ this is $(2+2)\cdot2^{1}=8$. A special solution can satisfy many polynomial relations, so a sufficiently rich dictionary is needed to reveal which relation actually identifies the equation rather than describing that particular solution.

\medskip\textbf{Step 7: why a closed form can fail to identify.} $u=e^{x+t}$ satisfies $u_t-u=0$, $u_x-u=0$, and consequences such as $u_{tx}-u=0$, $u_{xx}-u=0$. The solution alone does not say which relation is ``the'' governing equation: several equations explain the same special function. That is the intuition behind ``closed form $\Rightarrow$ does not identify''; the exact meaning of ``closed form'' in the theorem is the one fixed in the cited result, but the broad lesson is that an explicit solution may be too special to separate the true equation from competitors.

\medskip\textbf{The big picture.} The question changes from ``which linear differential operators annihilate $u$?'' to ``which polynomial differential relations does $u$ satisfy, and which genuinely identify its equation rather than merely describing a special solution?''. The dictionary makes the question finite-dimensional; $[F]$ accounts for all consequences of the proposed equation; $\mathcal J_u\cap\psi(V)$ detects extra relations. The crucial distinction is between \emph{fitting} a solution and \emph{identifying} its equation.

\subsection{Future comments: the linear and the nonlinear approach compared}\label{expl:E15}
\noindent\emph{(E15: Section 7.4, comment on the two approaches)}
You are not identifying the \emph{functions}; you are identifying the \emph{equations or relations} that the observed functions satisfy. Both approaches share this philosophy but use different mathematical objects.

\medskip\textbf{Linear case: look for an annihilating operator.} Given an observed $u(t,x)$ one searches for a constant-coefficient operator with $P(\partial_t,\partial_x)u=0$. For $u=e^{x+t}$: $(\partial_t-1)u=0$ and $(\partial_x-1)u=0$. One studies $\Ann(u)=\{P\in\C[\partial_t,\partial_x]:Pu=0\}$ and wants to characterize all the linear operators that annihilate $u$.

\medskip\textbf{Nonlinear case: search inside a finite dictionary.} For Burgers' $u_t+uu_x=0$ the product $uu_x$ is not a linear constant-coefficient operator applied to $u$. One chooses a finite dictionary, e.g.\ $\mathcal D=\{u_t,u_x,uu_x,u_{xx},u^2\}$, and searches for coefficients, not all zero, such that $c_1u_t+c_2u_x+c_3uu_x+c_4u_{xx}+c_5u^2=0$. This is a \emph{linear-algebra} problem in the unknown coefficients even though the candidate terms are nonlinear functions of $u$. If the data satisfy Burgers the coefficient vector $(1,0,1,0,0)$ works. The dictionary can be enlarged (higher-degree products, higher-order derivatives) but its limits must be decided in advance.

\medskip\textbf{What stays and what changes.}
\begin{center}\small
\begin{tabular}{@{}L{2.9cm}L{5.6cm}L{5.8cm}@{}}
\toprule
 & Linear approach & Nonlinear approach\\\midrule
Object observed & the function $u$ & the function $u$\\
Candidate relations & linear differential operators & differential polynomials\\
Search space & polynomial ring of constant-coefficient operators & a chosen finite dictionary of jet monomials\\
Main question & which operators annihilate $u$? & which candidate relations does $u$ satisfy?\\
Main limitation & structure of the annihilator & the dictionary may omit the true equation or contain too few informative terms\\
\bottomrule
\end{tabular}
\end{center}
The common principle is \emph{identifiability}: does the observed data contain enough information to distinguish the governing equation from alternative candidates?

\medskip\textbf{Where components and interactions enter.} For $u=u_1+u_2$ one can look at the components, the total, and an interaction term such as $\Eint=u-u_0-(u_1-u_0)-(u_2-u_0)$ (the additive form of Section~\ref{sec:inside}). \emph{Decomposition and interaction analysis} ask how the solution is assembled and how its components influence one another; \emph{equation identification} asks which differential relation governs the observed solution or its components. A linear annihilator may be used for a component when the assumptions allow, and a nonlinear dictionary to identify a nonlinear equation governing the total solution: complementary analyses, not automatically the same procedure.

\medskip\textbf{Key point.} The nonlinear dictionary approach generalizes the identification philosophy; it is not a direct replacement for the annihilator in every mathematical sense. The dictionary makes the search computationally manageable, but it also limits what can be discovered.

\subsection{The last section in plain words}\label{expl:E16}
\noindent\emph{(E16: Section 9, simple explanation of the last part)}
This section answers three simple questions: what can we \emph{never} do with this method, what do we \emph{not know yet}, and how much can we trust what is written.

\medskip\textbf{1. What is impossible.} These are proved facts, not gaps:
\begin{itemize}[nosep]
\item \emph{No finite proof of identification.} If $u$ is not an exponential-polynomial, no finite list of numbers (derivatives at a point, samples) can prove that $\Ann(u)=(P)$, because every finite list can also come from some exponential-polynomial. Example: you can refute ``$u$ is a sum of exponentials'' by finding a derivative that breaks it, but you can never confirm the opposite from finite data.
\item \emph{No splitting by terms.} Writing the solution as one piece per term of the operator fails: for the wave equation $\partial_t^2-c^2\partial_x^2$ only polynomial solutions split that way (Section~\ref{sec:groups}).
\item \emph{No stable removal of the interaction.} The interaction term cannot be switched off by a penalty or a regularization.
\item \emph{No identification with variable coefficients from one solution.}
\item \emph{No splitting on domains with a hole} (the logarithm example).
\item \emph{No finite product of single-term flows} reproduces non-commuting dynamics.
\end{itemize}

\medskip\textbf{2. What is open or conditional.} These are questions the paper cannot answer: KdV data with nonzero reflection coefficient; shock estimates uniform in time with several shock births; density of the sum of kernels on non-convex domains; exact values of an invariant $\delta_N(K)$; nonlinear asymptotics with derivatives; closedness of the sum of kernels for three or more differential factors. ``Open'' means unknown; ``conditional'' means true if an extra lemma holds.

\medskip\textbf{3. How reliable is it.} The paper is honest about this: many ingredients are classical (Nullstellensatz, Malgrange--Ehrenpreis, Cole--Hopf); the novelty is the \emph{package} of sharp statements, negative theorems and counterexamples, each with a status label. The reviewers were AI agents, not human referees; fifteen algebraic theorems were also machine-checked in Lean~4, while the analytic parts (distributions, Sobolev spaces, semigroups) were not, because the library is missing. Some citations were checked only on secondary sources, so before relying on a result one should trace it to its primary source.

\medskip\textbf{Take-home message.} Read the labels: \emph{proved} can be used, \emph{conditional} needs an extra lemma, \emph{open} is unknown, \emph{false} is a counterexample. The method is exact and has precise limits.
